\section{讨论}

负载规范决定哪些元件组合满足共模抵消要求。开路共模抵消固定两个系数，指定正电阻负载固定三个系数组合，两个不同正电阻负载则固定全部四个系数。元件恢复不等式进一步将这些系数要求转化为可直接检验的正值元件设计域。

P/Q 与 P/T 的区别给出了直接的设计步骤。对于开路目标，根据所需电压传递函数确定 a、b，再在式~\eqref{eq:8} 内选择 c、d。对于单一指定电阻负载，确定 b、c 和 a + gd，并利用剩余自由度构造支路。对于覆盖两个不同电阻负载的抵消要求，该拓扑必须采用相等的对应元件值。

\citet{casas1996basics} 在另一种电路配置中通过浮置耦合实现共模抵消。本文的对地梯形网络保留了流向地的电流路径，式~\eqref{eq:6} 显式给出了端接负载的作用。\citet{hornero2021cm} 通过共模响应描述了随频率变化的失配现象。本文的系数表示将完整的二阶支路响应与元件值、输出阻抗联系起来。

元件恢复结论要求元件值严格为正。当 \(q_1\) = 0 时，两个串联电阻合并为一个电阻；元件组合 (1,2,0,3) 和 (3/2,3/2,0,3) 得到相同的系数 (9,0,3,0)，式~\eqref{eq:7} 此时出现奇异。因此，双负载唯一性结论适用于第~\ref{sec:model-plane} 节定义的全正值二阶梯形网络。

对于导纳为 \(Y_L(s)\) 的频率相关线性负载，支路传递函数为 \(H_0/[1+Z_0Y_L(s)]\)。设计面向此类负载时，可先代入其导纳，再比较两个支路的响应。式~\eqref{eq:12} 则描述两个输出节点之间的耦合情形。


\subsection{设计与验证步骤}

滤波器规范可分四步转化为元件约束。首先，确定信号源平面、输出测量平面，以及各输出端所见的负载导纳。其次，选择该端接下所需的分母系数，并在正值元件域内选定剩余自由系数。再次，按式~\eqref{eq:7} 恢复元件值，并逐项检验式~\eqref{eq:8} 的不等式。最后，分别扰动元件值和两侧负载导纳，计算相应的 $G_{DC}$ 与 $G_{DD}$。精确等式确定标称设计；式~\eqref{eq:sensitivity} 和式~\eqref{eq:load-asymmetry} 则指出需要控制容差的组合。

双负载结论同时提供了一种验证思路。在两个已知电阻负载下测量完整支路响应，可约束四个系数。实验中需根据复频率响应数据估计这些系数，并给出不确定度区间。精确定理确定可辨识的模型结构；两个相近元件组合能够被区分到何种程度，则取决于测量不确定度与所选频点。

\subsection{待完成的实验验证}

\begin{quote}
\textbf{待补实验数据。} 当前结果包括精确代数推导与电路仿真。下一版预留 P/Q、P/T 在开路、10~k$\Omega$ 和 5~k$\Omega$ 端接下的实物电路比较。元件实测值、信号源与探头的复阻抗、频率响应、不确定度预算和重复性记录尚未采集。
\end{quote}

实验将采用相同的四元件梯形网络，以及表~\ref{tab:branches} 的元件比例。先分别测量各支路响应，再对支路对施加共模和差模激励。记录电阻、电容实际值后，可分别比较标称元件预测与实测元件预测。负载和探头导纳按输出测量平面接入带载模型。对于接近仪器噪声底的共模泄漏，将报告带界限的测量结果及相应检出限。

安装跨接电容后将进行第二轮采集，比较两种复增益及其比值。将信号源和探头配置表征为两个独立的线性支路等效网络，利用其带载传递函数和输出阻抗检验式~\eqref{eq:12} 的共同因子预测。记录的响应将与实测元件模型预测进行比较，并结合两者的不确定度判断符合程度。
